\section{Introduction}
\label{sec:intro}

Potentiostats built around Analog Devices' AD5940/AD5941 electrochemical
analog front end (AFE) have become a recognizable instrument class within
the open-source, low-cost hardware literature. The AD5940/AD5941 integrates
a potentiostat loop, a transimpedance amplifier with a selectable feedback
resistor ladder, a 16-bit sigma-delta ADC, and a hardware DFT accelerator on
a single part~\cite{devices2024ad5940}, which is why it recurs across
independently developed instruments: FreiStat~\cite{bill2023electrochemical},
HELPStat~\cite{bautista2025helpstat}, GaneStat~\cite{anshori2025ganestat},
and Vamos et al.'s HunStat2~\cite{vamos2025hunstat2} all use it as their AFE.
The instrument studied in this paper is our own AD5941/RP2040 potentiostat,
also named HunStat2, built directly from the AD5940 datasheet, standard
three-electrode potentiostat design practice, and the AD5940 evaluation-kit
reference firmware~\cite{ad5940_eis_github} (Section~\ref{sec:instrument-design}
details the components and signal-chain layout); it is listed separately
from Vamos et al.'s instrument of the same name in
Table~\ref{tab:related} to keep the two unambiguous. This is a meaningful
design choice, not an incidental one.
A discrete op-amp potentiostat -- the lineage that CheapStat and its
successors established~\cite{rowe2011cheapstat} -- puts the burden of
transimpedance gain accuracy, switch-matrix leakage, and reference-voltage
stability on the builder's own PCB layout and component tolerances. A
commercial AFE moves that burden onto a part that Analog Devices
manufactures, tests, and characterizes at volume; what a AD5940-based build
can still get wrong is confined to firmware -- how the chip's registers are
sequenced, how its raw codes are converted to physical units, and how that
logic is organized in code. That confinement is the premise of this paper:
if the analog front end is already commercial-grade, the object most worth
scrutinizing for reliability is the software layer sitting on top of it, not
the analog design.

The most recent survey of this instrument class, published in 2026,
substantiates the premise but also exposes a gap. Aviha and Slaughter's
review of modern potentiostat architectures explicitly covers ``Control
Software and User Interfaces'' as a section, citing FreiStat and Vamos et
al.'s HunStat2 by name~\cite{aviha2026modern}. Its treatment is entirely functional --
parameter configuration, real-time visualization, scripted automation -- and
at no point addresses how that software is structured: no discussion of
modularity, code duplication, encapsulation, or any other software-engineering
quality attribute appears anywhere in the review. Von Zuben et al.'s 2024
review of low-cost open-source potentiostats covers the electronics side of
DIY builds in comparable depth and likewise stops at the hardware/firmware
boundary~\cite{vonzuben2024lowcost}. Neither review, nor any of the
individual instrument papers surveyed in
Table~\ref{tab:related}, treats the firmware's software architecture as a
variable that affects instrument reliability. This is the gap this paper
addresses: we ask, for one concrete instrument, whether firmware software
quality can be measured, whether it can be improved with ordinary
refactoring technique, and whether that improvement can be shown -- not
argued -- to leave the instrument's actual measurements unchanged.

\begin{table}[h]
\centering
\caption{Software-architecture treatment in AD5940-lineage and comparable
low-cost potentiostat literature. ``Vamos et al.'' denotes the independently
published instrument of~\cite{vamos2025hunstat2}, which happens to share a
name with, but is a separate build from, the HunStat2 instrument audited in
this paper (``This work'').}
\label{tab:related}
\begin{threeparttable}
\begin{tabular}{@{}lllc@{}}
\toprule
Work & AFE & Methods & Firmware architecture assessed \\
\midrule
CheapStat~\cite{rowe2011cheapstat}       & Discrete op-amp & CV, SWV, LSV & No \\
FreiStat~\cite{bill2023electrochemical}  & AD5941      & CA, CV, DPV, EIS & No \\
HELPStat~\cite{bautista2025helpstat}     & AD5940      & CV, EIS      & No \\
GaneStat~\cite{anshori2025ganestat}      & Discrete op-amp & CV, CA, EIS & No \\
Vamos et al.~\cite{vamos2025hunstat2}    & AD5941      & CV, OCP, EIS & No \\
Review (2024)~\cite{vonzuben2024lowcost} & --          & --           & No (electronics only) \\
Review (2026)~\cite{aviha2026modern}     & --          & --           & No (functional only) \\
This work                                & AD5941      & CV, CA, SWV, DPV & \textbf{Yes} \\
\bottomrule
\end{tabular}
\end{threeparttable}
\end{table}

We audit our own HunStat2 firmware because it is the same instrument
independently validated at the metrology level in a companion
manuscript~\cite{clement2026linearity}, which found and corrected a
firmware calibration defect using a traceable dummy-cell protocol. That
work established \emph{that} the instrument's measurements can be trusted
once corrected; it did not examine \emph{how} the firmware that produces
those measurements is built, nor whether its structure poses a risk to that
correctness going forward. This paper is deliberately scoped as the
complementary half of that question. Where the companion manuscript treats
the dummy cell as the object of interest and the firmware as a black box to
be calibrated, this paper treats the firmware's internal structure as the
object of interest and reuses the same dummy-cell protocol only as a
correctness check -- to confirm that improving the software does not
regress the measurements the companion manuscript worked to validate.

Our contributions are:
\begin{enumerate}
\item A file-by-file object-oriented-design audit of the HunStat2 firmware,
  identifying (i) a measurement-configuration routine duplicated verbatim
  across three of five electrochemical-method classes, (ii) two entire
  subsystems -- a command-line parser and a second implementation of the
  electrochemical methods -- that are fully unreachable at runtime, and
  (iii) a live status-indicator bug that silently discards its input.
\item A scoped refactor that collapses (i) into one shared base class,
  removes (ii) after confirming via call-site tracing that it is dead, and
  fixes (iii), implemented as changes a reviewer can verify line-by-line
  against the original.
\item Hardware-measured validation of the refactor: the same CV, CA, SWV,
  and DPV sweeps against the same real dummy cell, run on the instrument
  before and after the refactor, quantifying how close ``after'' stays to
  ``before'' rather than merely asserting it.
\item A second, independent real load (a plain resistor, distinct from the
  dummy cell used for the primary comparison) run on the refactored
  firmware only, to corroborate that the post-refactor behavior generalizes
  beyond one specific load rather than having been tuned to it.
\end{enumerate}
